\chapter{Đứa Học Giỏi Chưa Chắc Yên}
\markboth{Đứa Học Giỏi Chưa Chắc Yên}{The Strong Student Is Not Necessarily at Peace}

\section{Điểm Cao Nhưng Mắt Rất Mệt}
\begin{truyen}{Điểm Cao Nhưng Mắt Rất Mệt}{High Scores but Tired Eyes}
\chuhoa{C}{ó em bài nào cũng tốt, điểm nào cũng cao, nhưng ánh mắt lúc nào cũng như đang nợ một thứ gì đó.}
\textit{(Some students do well on every test, yet their eyes always look like they owe something.)}

Các em không được phép yếu, không được phép trượt, không được phép thở ra lâu.
\textit{(They do not feel allowed to be weak, to slip, or to breathe for long.)}

Nhiều người chỉ nhìn thành tích và quên hỏi em sống có mệt không.
\textit{(Many people only look at performance and forget to ask whether the student is exhausted.)}

Học giỏi mà không yên thì cái giá cũng không nhỏ.
\textit{(Being excellent without peace comes with a serious cost.)}

\end{truyen}

\begin{baihoc}
Thành tích đẹp không có nghĩa là một đứa trẻ đang ổn.
\textit{(Beautiful performance does not guarantee that a child is okay.)}
\end{baihoc}

\begin{ghinhoanh}
\textbf{Short English to remember:}

Beautiful performance does not guarantee that a child is okay.

\end{ghinhoanh}

\begin{tuhocnhanh}
\textbf{Gợi ý để dạy và tự nghĩ / Teaching and reflection prompts}

1. Thành tích đang che điều gì? \textit{What is achievement hiding?}

2. Người dạy nên hỏi thêm điều gì? \textit{What should the teacher ask beyond scores?}

3. Học sinh giỏi cần được đỡ ở điểm nào? \textit{Where do strong students need support?}
\end{tuhocnhanh}
\ngancach

\section{Đứng Nhất Quá Lâu Nên Sợ Sai}
\begin{truyen}{Đứng Nhất Quá Lâu Nên Sợ Sai}{Staying on Top Too Long Can Create Fear of Failure}
\chuhoa{C}{ó em đã quen là người làm được nhanh nhất, đúng nhất, được gọi tên nhiều nhất.}
\textit{(Some students are used to being the fastest, most accurate, and most frequently praised.)}

Đến một lúc, các em không còn học chỉ để hiểu nữa, mà học để đừng rơi xuống.
\textit{(At some point, they no longer learn mainly to understand, but to avoid falling.)}

Nỗi sợ mất vị trí làm chuyện học bớt vui đi rất nhiều.
\textit{(The fear of losing position can drain much of the joy from learning.)}

Một đứa trẻ giỏi cũng có thể bị nhốt trong chính hình ảnh giỏi của mình.
\textit{(A strong student can become trapped inside the image of being strong.)}

\end{truyen}

\begin{baihoc}
Khen một đứa trẻ vì cố gắng thường cứu nó tốt hơn chỉ khen vì đứng đầu.
\textit{(Praising effort often protects a child better than praising rank alone.)}
\end{baihoc}

\begin{ghinhoanh}
\textbf{Short English to remember:}

Praising effort can protect a child better than praising rank alone.

\end{ghinhoanh}

\begin{tuhocnhanh}
\textbf{Gợi ý để dạy và tự nghĩ / Teaching and reflection prompts}

1. Em này đang học vì hiểu hay vì sợ rơi? \textit{Is this student learning to understand or to avoid falling?}

2. Kiểu khen nào đang tạo áp lực? \textit{What kind of praise is creating pressure?}

3. Làm sao để giữ động lực mà bớt sợ? \textit{How can motivation remain while fear decreases?}
\end{tuhocnhanh}
\ngancach

\section{Học Giỏi Nhưng Rất Cô Đơn}
\begin{truyen}{Học Giỏi Nhưng Rất Cô Đơn}{Excellent but Lonely}
\chuhoa{C}{ó em làm bài rất tốt nhưng giờ ra chơi lại ngồi một mình.}
\textit{(Some students perform very well yet sit alone during breaks.)}

Các em không gây chuyện, không làm phiền ai, nên sự cô đơn đó đi qua mắt người lớn rất nhẹ.
\textit{(They cause no trouble and disturb no one, so their loneliness passes softly through adult attention.)}

Người ta thường thương đứa ồn hơn đứa lặng.
\textit{(People often notice the loud child sooner than the quiet one.)}

Nhưng có những nỗi buồn đi rất êm trong một đứa trẻ học giỏi.
\textit{(Yet some sadness moves very quietly inside high-performing students.)}

\end{truyen}

\begin{baihoc}
Một đứa trẻ ngoan và giỏi vẫn có thể đang thiếu một chỗ để thuộc về.
\textit{(A good and capable child may still be lacking a place to belong.)}
\end{baihoc}

\begin{ghinhoanh}
\textbf{Short English to remember:}

A capable child may still be lacking a place to belong.

\end{ghinhoanh}

\begin{tuhocnhanh}
\textbf{Gợi ý để dạy và tự nghĩ / Teaching and reflection prompts}

1. Điều gì dễ bị bỏ sót ở một em học giỏi? \textit{What is easily overlooked in a strong student?}

2. Vì sao nỗi cô đơn này khó thấy? \textit{Why is this loneliness hard to notice?}

3. Lớp học có thể ấm hơn bằng cách nào? \textit{How can the classroom become warmer?}
\end{tuhocnhanh}
\ngancach

\section{Được Khen Nhiều Nên Không Dám Yếu}
\begin{truyen}{Được Khen Nhiều Nên Không Dám Yếu}{Praised So Much That Weakness Feels Forbidden}
\chuhoa{N}{hiều em càng được tin, càng được khen, lại càng khó nói ra lúc mình mệt.}
\textit{(The more some students are praised and trusted, the harder it becomes for them to admit exhaustion.)}

Các em sợ phụ lòng thầy cô, sợ làm cha mẹ hụt hẫng, sợ hình ảnh mình trong mắt người khác thay đổi.
\textit{(They fear disappointing teachers, parents, and the image others hold of them.)}

Nên có em vẫn cười, vẫn làm, vẫn nộp đủ, nhưng bên trong thì mòn dần.
\textit{(So they keep smiling, working, and submitting, while quietly wearing down inside.)}

Đôi khi học sinh giỏi cũng cần được phép yếu như một người bình thường.
\textit{(Sometimes strong students need permission to be weak like any normal person.)}

\end{truyen}

\begin{baihoc}
Kỳ vọng đẹp nếu thiếu chỗ nghỉ sẽ thành gánh nặng đẹp.
\textit{(Beautiful expectations become beautifully packaged burdens when there is no place to rest.)}
\end{baihoc}

\begin{ghinhoanh}
\textbf{Short English to remember:}

Beautiful expectations can become burdens when there is no place to rest.

\end{ghinhoanh}

\begin{tuhocnhanh}
\textbf{Gợi ý để dạy và tự nghĩ / Teaching and reflection prompts}

1. Vì sao học sinh giỏi hay khó nói ra sự mệt? \textit{Why do strong students struggle to admit exhaustion?}

2. Kỳ vọng đẹp biến thành gánh ở chỗ nào? \textit{Where does healthy expectation become burden?}

3. Người dạy có thể cho phép điều gì nhiều hơn? \textit{What permission can the teacher give more often?}
\end{tuhocnhanh}
\ngancach

\section{Được Điểm Chín Mà Muốn Khóc}
\begin{truyen}{Được Điểm Chín Mà Muốn Khóc}{Getting a Nine and Wanting to Cry}
\chuhoa{C}{ó em cầm bài 9 điểm mà mặt tái đi.}
\textit{(Some students receive a 9 out of 10 and go pale.)}

Tôi hỏi: “Bài tốt mà, sao vậy?”
\textit{(I asked, “It’s a good paper. What happened?”)}

Em nói nhỏ: “Mẹ em sẽ hỏi câu sai nằm ở đâu.”
\textit{(The student whispered, “My mother will ask where the missing point went.”)}

Lúc đó mới thấy có những đứa trẻ không sợ điểm thấp. Chúng sợ chưa đủ hoàn hảo.
\textit{(That was when I saw that some children do not fear low scores. They fear being less than perfect.)}

\end{truyen}

\begin{baihoc}
Một đứa trẻ luôn phải đủ mười rất khó thấy nhẹ trong chuyện học.
\textit{(A child who always has to be perfect will rarely feel light inside learning.)}
\end{baihoc}

\begin{ghinhoanh}
\textbf{Short English to remember:}

A child who always has to be perfect will rarely feel light inside learning.

\end{ghinhoanh}

\begin{tuhocnhanh}
\textbf{Gợi ý để dạy và tự nghĩ / Teaching and reflection prompts}

1. Nỗi sợ thật của em này là gì? \textit{What is this student truly afraid of?}

2. Vì sao điểm 9 vẫn có thể làm đau? \textit{Why can a 9 still hurt?}

3. Người dạy có thể giúp hạ chuẩn hoàn hảo thế nào? \textit{How can a teacher soften perfectionism?}
\end{tuhocnhanh}
\ngancach

\section{Cả Lớp Nhờ, Một Em Mệt}
\begin{truyen}{Cả Lớp Nhờ, Một Em Mệt}{The Whole Class Leans on One Tired Student}
\chuhoa{C}{ó em học giỏi nên ai cũng nhờ: hỏi bài, mượn vở, xin chỉ cách làm.}
\textit{(Some students do well, so everyone leans on them: asking for help, borrowing notes, requesting explanations.)}

Bề ngoài đó là sự quý mến.
\textit{(On the surface, it looks like affection.)}

Nhưng nếu kéo dài, nó làm một học sinh giỏi thấy mình luôn phải đủ, luôn phải biết, luôn phải sẵn sàng.
\textit{(But over time, it makes a strong student feel they must always know, always be ready, always be enough.)}

Có những đứa trẻ mệt vì quá lâu không được phép yếu trước người khác.
\textit{(Some children grow exhausted because they have not been allowed to look weak in front of others.)}

\end{truyen}

\begin{baihoc}
Được tin cậy là tốt, nhưng bị xem như lúc nào cũng phải gánh là một áp lực khác.
\textit{(Being trusted is good, but being treated as someone who must always carry others is another pressure.)}
\end{baihoc}

\begin{ghinhoanh}
\textbf{Short English to remember:}

Being trusted is good, but always having to carry others is another pressure.

\end{ghinhoanh}

\begin{tuhocnhanh}
\textbf{Gợi ý để dạy và tự nghĩ / Teaching and reflection prompts}

1. Vì sao học sinh giỏi cũng có thể kiệt? \textit{Why can strong students burn out too?}

2. Điều gì nhìn giống tin cậy nhưng hóa ra là gánh? \textit{What looks like trust but becomes burden?}

3. Mình có đang vô tình đặt quá nhiều lên một em nào không? \textit{Are we placing too much on one student?}
\end{tuhocnhanh}
\ngancach

\section{Giỏi Nên Cũng Bị Ghen}
\begin{truyen}{Giỏi Nên Cũng Bị Ghen}{Excellence Can Also Attract Jealousy}
\chuhoa{C}{ó em học tốt, cư xử đàng hoàng, nhưng vẫn bị nói xấu chỉ vì luôn nổi lên.}
\textit{(Some students do well and behave properly, yet are still talked about badly simply because they stand out.)}

Tuổi học trò không chỉ có quý mến. Nó cũng có ganh tị rất non mà rất thật.
\textit{(School age does not contain only affection. It also contains jealousy that is immature yet very real.)}

Nhiều em học giỏi bị cô đơn không phải vì kiêu, mà vì bị để ý theo cách không dễ chịu.
\textit{(Many strong students become lonely not because they are proud, but because attention toward them turns uncomfortable.)}

Giỏi đôi khi cũng là một kiểu dễ bị nhắm vào.
\textit{(Being strong can sometimes make a child easier to target.)}

\end{truyen}

\begin{baihoc}
Một học sinh giỏi cũng cần được bảo vệ khỏi sự ghen ghét của tập thể.
\textit{(A high-performing student also needs protection from group jealousy.)}
\end{baihoc}

\begin{ghinhoanh}
\textbf{Short English to remember:}

A high-performing student also needs protection from group jealousy.

\end{ghinhoanh}

\begin{tuhocnhanh}
\textbf{Gợi ý để dạy và tự nghĩ / Teaching and reflection prompts}

1. Vì sao người nổi lên dễ bị nhắm tới? \textit{Why are standout students easy targets?}

2. Ganh tị tuổi học trò biểu hiện ra sao? \textit{How does school-age jealousy show up?}

3. Người dạy cần giữ công bằng thế nào? \textit{How should a teacher protect fairness?}
\end{tuhocnhanh}
\ngancach

\section{Cho Em Được Nghỉ Làm Người Giỏi Một Hôm}
\begin{truyen}{Cho Em Được Nghỉ Làm Người Giỏi Một Hôm}{Let the Student Rest from Being the “Strong One”}
\chuhoa{C}{ó hôm tôi nói với một em học giỏi: “Hôm nay em không cần trả lời nhiều. Em cứ ngồi học bình thường thôi.”}
\textit{(One day I told a strong student, “Today you do not have to answer much. Just sit and learn like everyone else.”)}

Em nhìn tôi như không quen với việc đó.
\textit{(The student looked at me as if unfamiliar with that permission.)}

Có những đứa trẻ giỏi đến mức quên mất cảm giác được phép không gồng.
\textit{(Some excellent children become so used to performing that they forget what it feels like not to tense up.)}

Được cho nghỉ khỏi vai giỏi đôi khi cũng là một sự cứu.
\textit{(Being allowed to rest from the role of the “good one” can itself be a rescue.)}

\end{truyen}

\begin{baihoc}
Học sinh giỏi cũng cần những khoảng không phải chứng minh gì cả.
\textit{(Strong students also need spaces where they do not have to prove anything.)}
\end{baihoc}

\begin{ghinhoanh}
\textbf{Short English to remember:}

Strong students also need spaces where they do not have to prove anything.

\end{ghinhoanh}

\begin{tuhocnhanh}
\textbf{Gợi ý để dạy và tự nghĩ / Teaching and reflection prompts}

1. Vì sao “vai người giỏi” cũng làm mệt? \textit{Why can the role of the strong student become exhausting?}

2. Người dạy có thể cho các em nghỉ ở đâu? \textit{Where can teachers let these students rest?}

3. Mình có đang dùng các em quá nhiều không? \textit{Are we relying on them too much?}
\end{tuhocnhanh}
